% !TeX root=../main.tex
\chapter{جمع‌بندی و کارهای آینده}\label{Chap:Conclusion}

\section{جمع‌بندی پروژه}

رشد ارتباط و واژگان در سه سال نخست زندگی با تعامل مداوم کودک و اطرافیان شکل می‌گیرد. والد در فعالیت‌های روزانه فرصت‌های فراوانی برای نامیدن اشیا، پاسخ‌دادن به اشاره و آوا، تکرار واژه و ایجاد گفت‌وگو دارد؛ بااین‌حال، فراهم‌کردن تمرین‌های منظم، متنوع و متناسب با توانایی کودک همیشه آسان نیست. پروژهٔ «نوواژه» با هدف ایجاد یک ابزار همراه برای سامان‌دادن این تمرین‌ها و افزایش مشارکت والد و کودک طراحی و پیاده‌سازی شد.

در نوواژه، واژه‌ها در دسته‌هایی مانند میوه‌ها، حیوانات، غذاها، رنگ‌ها و وسایل نقلیه سازمان‌دهی شده‌اند. هر واژه با تصویر و نمونهٔ صوتی ارائه می‌شود و کودک می‌تواند پس از دیدن تصویر، تلفظ آن را بشنود. در بخش آزمون، واژه پخش می‌شود و کودک باید تصویر متناظر را از میان گزینه‌ها انتخاب کند. این روش، پیوند میان شکل شنیداری واژه و مرجع دیداری آن را تمرین می‌دهد. ثبت پاسخ‌ها نیز امکان مشاهدهٔ روند فعالیت و اولویت‌دادن به واژه‌های نیازمند تمرین بیشتر را فراهم می‌کند.

برای حفظ مشارکت کودک، بازی‌هایی مانند حافظه، مرتب‌سازی رنگ‌ها، تشخیص گزینهٔ متفاوت، تطبیق سایه و چهره‌سازی در برنامه قرار گرفته‌اند. این بازی‌ها در کنار آموزش مستقیم واژه، توجه دیداری، یادآوری، دسته‌بندی و تشخیص شباهت و تفاوت را درگیر می‌کنند. بازخورد کوتاه، دکمه‌های بزرگ و تعداد محدود گزینه‌ها نیز با توجه به گروه سنی کودک انتخاب شده‌اند.

بخش والدین یکی از اجزای مهم سامانه است. والد می‌تواند حساب کاربری ایجاد کند، برای یک یا دو کودک نمایه بسازد، کودک فعال را انتخاب کند و زمان مجاز فعالیت را تنظیم نماید. همچنین امکان افزودن دسته و واژهٔ شخصی با تصویر دلخواه و صدای ضبط‌شده وجود دارد. این قابلیت اجازه می‌دهد محتوای برنامه با محیط واقعی کودک، اعضای خانواده، اسباب‌بازی‌ها و اشیای روزمره هماهنگ شود. حالت کودک و رمز چهاررقمی نیز برای کاهش خروج ناخواسته از برنامه در نظر گرفته شده‌اند.

از دید فنی، نوواژه به‌صورت یک سامانهٔ چندسکویی توسعه یافته است. رابط مشترک اندروید و آی‌اواس با کاتلین چندسکویی و کامپوز چندسکویی پیاده‌سازی شده و کارساز کتور وظایفی مانند احراز هویت، مدیریت محتوا، تولید آزمون و نگهداری پاسخ‌ها را انجام می‌دهد. جداسازی برنامهٔ موبایل، منطق مشترک و کارساز، توسعه و نگهداری قابلیت‌ها را ساده‌تر کرده و زمینهٔ گسترش آیندهٔ سامانه را فراهم آورده است.

حاصل پروژه، نمونه‌ای کاربردی است که آموزش دیداری و شنیداری، بازی، شخصی‌سازی محتوا، مدیریت والد و ثبت پیشرفت را در یک برنامه یکپارچه می‌کند. با وجود این، نوواژه جایگزین گفتاردرمانی یا ارزیابی تخصصی نیست. نتیجهٔ آزمون‌ها تنها عملکرد کودک در فعالیت‌های درون برنامه را نشان می‌دهد و برای قضاوت دربارهٔ تلفظ، شنوایی، زبان بیانی یا وضعیت رشدی کافی نیست. ارزش اصلی برنامه در ایجاد فرصت‌های کوتاه، ساخت‌یافته و قابل تکرار برای تعامل والد و کودک است.

\section{محدودیت‌های نسخهٔ فعلی}

نسخهٔ فعلی همهٔ بازی‌ها را با ساختاری نسبتاً ثابت ارائه می‌کند. تعداد گزینه‌ها، پیچیدگی پرسش و قواعد هر بازی برای همهٔ کودکان تقریباً یکسان‌اند؛ در نتیجه ممکن است یک فعالیت برای کودک تازه‌کار دشوار و برای کودک باتجربه تکراری باشد. همچنین بازی‌ها به مرحله‌های پیوسته تقسیم نشده‌اند و کودک پس از پایان یک دور، مسیر روشنی برای حرکت از فعالیت ساده به پیچیده مشاهده نمی‌کند.

وضعیت آموخته‌شدن واژه نیز بر تعداد پاسخ‌های درست تکیه دارد. این روش برای نسخهٔ اولیه ساده و قابل فهم است، اما عواملی مانند فاصلهٔ زمانی میان تمرین‌ها، سرعت پاسخ، تنوع تصویر و استفادهٔ واژه در زندگی واقعی را در نظر نمی‌گیرد. از سوی دیگر، ظاهر برنامه یک تم اصلی دارد و امکان انتخاب فضای بصری متناسب با علاقهٔ کودک هنوز فراهم نشده است.

ارزیابی نسخهٔ حاضر بیشتر بر درستی عملکرد نرم‌افزار و اجرای قابلیت‌ها متمرکز بوده است. برای بررسی اثر واقعی برنامه بر تعامل والد و کودک، سهولت استفاده و یادگیری واژه، مطالعه با حضور خانواده‌ها و متخصصان گفتار و زبان ضروری خواهد بود. همچنین برای انتشار عمومی باید آزمون‌های گسترده‌تر روی دستگاه‌های مختلف، سیاست دقیق نگهداری داده و زیرساخت عملیاتی امن در نظر گرفته شود.

\section{کارهای آینده}

\subsection{سطح‌بندی بازی‌ها}

مهم‌ترین توسعهٔ آینده، افزودن سطح دشواری به بازی‌هاست. هر بازی می‌تواند حداقل سه سطح «آسان»، «متوسط» و «سخت» داشته باشد. سطح مناسب در ابتدا توسط والد انتخاب می‌شود و در نسخه‌های بعد می‌تواند بر اساس عملکرد کودک به‌طور خودکار پیشنهاد گردد.

در سطح آسان، تعداد گزینه‌ها کم، تصویرها کاملاً متفاوت، زمان پاسخ آزاد و راهنمای دیداری بیشتر است. در سطح متوسط، تعداد گزینه‌ها افزایش می‌یابد و گزینه‌های نزدیک‌تر از نظر شکل یا گروه معنایی نمایش داده می‌شوند. در سطح سخت، سرنخ‌ها کمتر، تنوع محتوا بیشتر و فعالیت نیازمند دقت یا حافظهٔ بالاتری خواهد بود. برای نمونه:

\begin{itemize}
\item در بازی حافظه، سطح دشواری با تعداد جفت کارت‌ها و شباهت تصویرها تغییر می‌کند؛
\item در بازی رنگ‌ها، تعداد رنگ‌ها، سرعت یا تعداد اشیای قابل جابه‌جایی افزایش می‌یابد؛
\item در بازی سایه‌ها، گزینه‌ها از نظر شکل به یکدیگر نزدیک‌تر می‌شوند؛
\item در بازی «کدام متفاوت است؟»، تفاوت آشکار در سطح آسان و تفاوت معنایی دقیق‌تر در سطح سخت استفاده می‌شود؛
\item در آزمون شنیداری، تعداد گزینه‌ها و شباهت واژه‌ها به‌تدریج بیشتر می‌شود.
\end{itemize}

انتخاب سطح نباید تنها بر یک پاسخ اشتباه یا درست متکی باشد. بهتر است چند جلسهٔ اخیر، درصد پاسخ درست، تعداد تلاش و نیاز کودک به تکرار صوت بررسی شوند. امکان بازگشت به سطح ساده‌تر نیز باید همیشه وجود داشته باشد تا بازی باعث اضطراب یا رهاکردن فعالیت نشود. سطح هر بازی به‌صورت جداگانه برای هر کودک ذخیره خواهد شد؛ زیرا ممکن است کودک در حافظه قوی‌تر و در تشخیص رنگ یا واژه به تمرین بیشتری نیاز داشته باشد.

\subsection{افزودن مرحله به بازی‌ها}

در نسخهٔ آینده، هر بازی می‌تواند به مجموعه‌ای از مرحله‌های کوتاه تقسیم شود. هر مرحله هدف مشخصی دارد و پس از تکمیل آن، مرحلهٔ بعد باز می‌شود. برای مثال، مرحلهٔ نخست بازی حافظه با دو جفت کارت آغاز می‌شود، مرحلهٔ دوم سه جفت و مرحله‌های بعدی تعداد بیشتری کارت یا تصویرهای مشابه‌تر ارائه می‌کنند.

ساختار مرحله‌ای باید پیشرفت را برای کودک قابل مشاهده کند، بدون آنکه رقابت یا فشار نامناسب ایجاد شود. نمایش مسیر ساده، ستاره یا نشان کوتاه و پیام تشویقی می‌تواند حس پیشرفت ایجاد کند. در عین حال، والد باید بتواند مرحله‌ای را تکرار، مرحلهٔ دشوار را موقتاً رد یا بازی را از سطح مناسب آغاز کند.

برای هر مرحله می‌توان اطلاعات زیر را ذخیره کرد:

\begin{itemize}
\item وضعیت باز یا تکمیل‌شده؛
\item تعداد دفعات تلاش؛
\item بهترین نتیجه و آخرین نتیجه؛
\item زمان تقریبی انجام؛
\item واژه‌ها یا مهارت‌های تمرین‌شده؛
\item سطح پیشنهادی برای جلسهٔ بعد.
\end{itemize}

مرحله‌ها نباید فقط با افزایش تعداد گزینه‌ها دشوار شوند. تنوع تصویر، ترکیب دسته‌ها، کاهش تدریجی راهنما، فاصلهٔ زمانی میان نمایش و پاسخ و نیاز به بازیابی فعال نیز می‌توانند پیچیدگی را افزایش دهند. طراحی مرحله‌ای همچنین امکان افزودن محتوای تازه بدون تغییر ساختار اصلی بازی را فراهم می‌کند.

\subsection{افزودن تم دخترانه و پسرانه}

یکی دیگر از قابلیت‌های پیشنهادی، امکان انتخاب تم ظاهری دخترانه یا پسرانه است. این تم می‌تواند رنگ‌های رابط، شخصیت راهنما، تصویر پس‌زمینه، شکل دکمه‌ها و برخی پویانمایی‌های تشویقی را تغییر دهد. انتخاب در بخش والدین انجام می‌شود و برای هر کودک به‌صورت جداگانه ذخیره خواهد شد.

در کنار دو تم پیشنهادی، حفظ یک تم عمومی و خنثی نیز مناسب است تا والد بتواند صرفاً بر اساس علاقهٔ کودک انتخاب کند. تم نباید محتوای آموزشی، سطح دشواری یا دسترسی به بازی‌ها را محدود سازد و نباید کلیشه‌های جنسیتی را به کودک تحمیل کند. هدف از تم، افزایش حس تعلق و جذابیت بصری است، نه جداسازی توانایی‌ها یا فعالیت‌ها.

برای پیاده‌سازی مناسب، رنگ‌ها باید خوانایی کافی داشته باشند، تصویرها موجب شلوغی صفحه نشوند و صداها کوتاه و قابل کنترل باقی بمانند. همچنین بهتر است امکان انتخاب جداگانهٔ رنگ اصلی و شخصیت راهنما وجود داشته باشد. در این صورت، کودک می‌تواند ظاهر مورد علاقهٔ خود را بدون وابستگی کامل به یک عنوان از پیش تعیین‌شده انتخاب کند.

\subsection{ارزیابی و توسعه‌های تکمیلی}

پس از افزودن سطح، مرحله و تم، لازم است نسخهٔ جدید با والدین و کودکان در محیط کنترل‌شده آزمایش شود. شاخص‌هایی مانند مدت مشارکت، تعداد خطا، نیاز به کمک والد، سهولت جابه‌جایی میان صفحه‌ها و رضایت خانواده می‌توانند برای اصلاح طراحی استفاده شوند. دریافت نظر آسیب‌شناس گفتار و زبان نیز به مناسب‌بودن واژه‌ها، ترتیب دشواری و نوع بازخورد کمک می‌کند.

در ادامه می‌توان الگوریتم مرور واژه را بهبود داد تا زمان تمرین، پاسخ‌های اخیر و فاصلهٔ میان جلسات را در نظر بگیرد. پشتیبانی بهتر از کار آفلاین، همگام‌سازی امن داده، گزارش ساده‌تر برای والد، دسترس‌پذیری برای کودکان دارای نیازهای متفاوت و آزمایش روی دامنهٔ بیشتری از تلفن‌ها نیز از توسعه‌های تکمیلی محسوب می‌شوند.

\section{نتیجهٔ نهایی}

نوواژه بستری توسعه‌پذیر برای ترکیب آموزش واژه، بازی و مشارکت والد فراهم کرده است. قابلیت‌های موجود امکان استفادهٔ روزمره و شخصی‌سازی محتوا را ایجاد می‌کنند و معماری چندسکویی، افزودن قابلیت‌های آینده را ممکن می‌سازد. سطح‌بندی بازی‌ها، طراحی مرحله‌های پیوسته و ارائهٔ تم‌های قابل انتخاب می‌توانند تجربه را متناسب‌تر، جذاب‌تر و پایدارتر کنند. موفقیت این مسیر به طراحی متناسب با سن، همراهی والد، ارزیابی با کاربران واقعی و حفظ مرز روشن میان ابزار آموزشی و درمان تخصصی وابسته خواهد بود.
